\documentclass[aspectratio=169,11pt]{beamer}

\usepackage[spanish,es-nodecimaldot]{babel}
\usepackage[T1]{fontenc}
\usepackage[utf8]{inputenc}
\usepackage{xcolor}
\usepackage{booktabs}
\usepackage{listings}

\definecolor{azul}{HTML}{174A7E}
\definecolor{azulclaro}{HTML}{DCECF8}
\definecolor{verde}{HTML}{2E7D32}
\definecolor{gris}{HTML}{F3F5F7}
\definecolor{naranja}{HTML}{A64B00}
\usetheme{Madrid}
\setbeamercolor{structure}{fg=azul}
\setbeamercolor{frametitle}{fg=azul}
\setbeamercolor{block title}{bg=azul,fg=white}
\setbeamercolor{block body}{bg=azulclaro,fg=black!85}
\setbeamertemplate{navigation symbols}{}
\setbeamertemplate{footline}{\hfill\insertframenumber\hspace{0.4cm}\vspace{0.2cm}}

\lstdefinestyle{codigo}{
  language=Python,
  basicstyle=\ttfamily\scriptsize,
  backgroundcolor=\color{gris},
  frame=single,
  breaklines=true,
  keepspaces=true,
  showstringspaces=false,
  columns=fullflexible,
  keywordstyle=\color{azul}\bfseries,
  commentstyle=\color{verde},
  stringstyle=\color{naranja}
}

\title{Del CSV al modelo}
\subtitle{Analizar, detectar señales y predecir con datos de ventas}
\author{Laboratorio de Inteligencia Artificial}
\date{}

\begin{document}

\begin{frame}
  \titlepage
\end{frame}

\begin{frame}{La pregunta de hoy}
  \begin{center}
    \Large ¿Podemos estimar cuántas unidades se venderán\
    usando visitas y descuento?
  \end{center}
  \vfill
  \begin{itemize}
    \item Partimos de una tabla CSV, como en la primera práctica.
    \item Usamos limpieza y agrupación antes de modelar.
    \item Evaluamos una predicción con datos que el modelo no vio.
  \end{itemize}
\end{frame}

\begin{frame}{El recorrido: de tabla a evidencia}
  \begin{enumerate}
    \item Cargar y revisar: tipos, faltantes, rangos.
    \item Transformar: crear el importe y agrupar por ciudad.
    \item Investigar: marcar posibles valores atípicos.
    \item Aprender: ajustar una regresión a partir del CSV.
    \item Evaluar: comparar el error con una predicción de referencia.
  \end{enumerate}
  \vfill
  \begin{block}{Conexión con clases anteriores}
    Es aprendizaje supervisado: el CSV contiene entradas y una respuesta conocida para aprender, como en la predicción de estaturas.
  \end{block}
\end{frame}

\begin{frame}{El conjunto de datos}
  \begin{center}
    \begin{tabular}{lrrr}
      \toprule
      fecha & visitas & descuento\_pct & unidades \\
      \midrule
      2026-01-01 & 95 & 7 & 12 \\
      2026-01-02 & 372 & 25 & 39 \\
      2026-01-03 & 257 & 9 & 29 \\
      \bottomrule
    \end{tabular}
  \end{center}
  \begin{itemize}
    \item 240 días sintéticos; hay una celda faltante y dos ventas inusuales.
    \item La columna \texttt{unidades} es el objetivo que queremos estimar.
    \item Las variables \texttt{visitas} y \texttt{descuento\_pct} son las entradas.
  \end{itemize}
\end{frame}

\begin{frame}[fragile]{1. Leer y auditar antes de analizar}
  \begin{lstlisting}[style=codigo]
with open("ventas_clase.csv", newline="") as archivo:
    lector = csv.DictReader(archivo)
    filas = list(lector)

print(lector.fieldnames)
print(filas[:2])
  \end{lstlisting}
  \begin{block}{Por qué importa}
    CSV guarda texto. Las conversiones a número deben ser explícitas, y una celda vacía requiere una decisión antes de calcular.
  \end{block}
\end{frame}

\begin{frame}[fragile]{2. Limpiar, transformar y resumir}
  \begin{lstlisting}[style=codigo]
limpios = []
for fila in filas:
    try:
        unidades = float(fila["unidades"])
        precio = float(fila["precio_unitario"])
    except ValueError:
        continue
    fila["importe"] = unidades * precio
    limpios.append(fila)
  \end{lstlisting}
  \begin{itemize}
    \item La práctica descarta filas incompletas solo en los campos necesarios.
    \item El resumen por ciudad describe el CSV; no explica por qué difieren las ciudades.
  \end{itemize}
\end{frame}

\begin{frame}{3. Valores atípicos: señal para investigar}
  Para una variable numérica, el rango intercuartílico es $IQR=Q_3-Q_1$.
  \[
    \text{marcar si } x < Q_1-1.5\,IQR
    \quad\text{o}\quad x > Q_3+1.5\,IQR
  \]
  \begin{itemize}
    \item La regla detecta valores alejados del centro de la distribución.
    \item Una promoción exitosa puede ser un valor extremo real.
    \item Igual que con Benford: anomalía no significa fraude; hay que revisar contexto y origen.
  \end{itemize}
\end{frame}

\begin{frame}{4. Entrenamiento y prueba en orden temporal}
  \begin{center}
    \begin{tabular}{c|c}
      \toprule
      \textbf{Primeros días: entrenamiento} & \textbf{Últimos días: prueba} \\
      \midrule
      aprende los coeficientes & mide errores en días no usados \\
      \bottomrule
    \end{tabular}
  \end{center}
  \vfill
  \begin{itemize}
    \item En datos temporales, no mezclamos el futuro con el pasado.
    \item Un buen resultado de entrenamiento no basta: interesa la prueba.
    \item No usamos columnas que revelen directamente la respuesta: eso sería fuga de información.
  \end{itemize}
\end{frame}

\begin{frame}{5. Un modelo pequeño y transparente}
  Suponemos una relación lineal:
  \[
    \widehat{unidades}=b_0+b_1\,visitas+b_2\,descuento
  \]
  \begin{itemize}
    \item \texttt{numpy.linalg.lstsq} busca coeficientes que ajusten el entrenamiento.
    \item No es una red neuronal, pero practica el mismo ciclo: ejemplos, ajuste y evaluación.
    \item Los coeficientes describen asociaciones en estos datos sintéticos, no causalidad.
  \end{itemize}
\end{frame}

\begin{frame}{¿Qué significa el error MAE?}
  \[
    MAE=\frac{1}{n}\sum_{i=1}^{n}|y_i-\hat{y}_i|
  \]
  \begin{itemize}
    \item Si el MAE es 4, la predicción se equivoca en promedio por 4 unidades.
    \item Comparamos con una referencia simple: predecir siempre la media del entrenamiento.
    \item El modelo solo es útil si supera una referencia razonable en datos de prueba.
  \end{itemize}
\end{frame}

\begin{frame}[fragile]{El núcleo del ajuste en pocas líneas}
  \begin{lstlisting}[style=codigo]
x_train = np.column_stack([np.ones(len(train)), train_x])
x_test = np.column_stack([np.ones(len(test)), test_x])
b = np.linalg.lstsq(x_train, y_train, rcond=None)[0]
prediccion = x_test @ b
mae = np.mean(np.abs(y_test - prediccion))
  \end{lstlisting}
  \begin{block}{Antes de ejecutar}
    ¿Qué variable creen que pesará más? ¿Qué cambiaría si quitamos el descuento de las entradas?
  \end{block}
\end{frame}

\begin{frame}{Resultados: describir, no exagerar}
  \begin{itemize}
    \item Revisar cuántas filas se conservaron y cuántos posibles atípicos aparecieron.
    \item Comparar MAE de referencia y MAE del modelo.
    \item Si falla, investigar: ¿datos ruidosos, relación no lineal, cambios con el tiempo?
  \end{itemize}
  \vfill
  \begin{alertblock}{Límite del ejemplo}
    Los datos son sintéticos y sencillos. Una asociación predictiva no demuestra que aumentar visitas o descuentos cause más ventas.
  \end{alertblock}
\end{frame}

\begin{frame}{Conexión con el robot y otros temas}
  \begin{itemize}
    \item Un CSV de episodios del robot puede tener sensores, acción, choque y recompensa.
    \item El mismo flujo permite predecir una recompensa o detectar episodios extraños.
    \item Embeddings y RAG son otra ruta: buscar semánticamente en texto asociado a los registros, no sustituir esta auditoría tabular.
  \end{itemize}
  \vfill
  \begin{block}{Idea común}
    Los modelos dependen de la calidad de los datos, de una pregunta clara y de una evaluación honesta.
  \end{block}
\end{frame}

\begin{frame}{Práctica: cambia una cosa y comprueba}
  \begin{enumerate}
    \item Ejecutar \texttt{python3 generar\_ventas.py}.
    \item Ejecutar \texttt{python3 analizar\_ventas.py}.
    \item Cambiar las entradas a solo \texttt{visitas}; comparar los MAE.
    \item Probar otra proporción temporal, por ejemplo 70/30.
    \item Inspeccionar las fechas marcadas como atípicas antes de decidir eliminarlas.
  \end{enumerate}
  \vfill
  \textbf{Entrega:} cambio, predicción previa, resultado observado y una limitación del análisis.
\end{frame}

\begin{frame}{Cierre}
  \begin{center}
    \Large CSV $\rightarrow$ preguntas $\rightarrow$ modelo $\rightarrow$ evaluación
  \end{center}
  \vfill
  \centering
  Un modelo no arregla datos confusos ni convierte una correlación en causa.
\end{frame}

\end{document}